% !TeX root = ../papers/nc-p3.tex
\section{反対階数八の分岐を中央二次元の中心で検出する}
\label{sec:branch-middle-defect-bound}
\begin{theorem}\label{thm:branch-middle-defect-bound}
曲面射影のない標準正則窓に片側分岐があるとする。反対階数8ならどちらかの向きに
中央支持2の分岐中心があり、その三テンソルは中央鉛筆 $L_2$ 型である。
反対階数9でこのような中心がないなら、全順・逆中心の中央支持は3であり、
各後続直線（逆向きでは前点直線）の一般点の中央支持は4である。
ここでは純関係を許す。
\end{theorem}
\begin{proof}
中心 $p$ の中央支持を $m$ とし、順向きとして位置0に置く。
後続直線 $L$ の一般点 $x$ の両端支持は3なので中央支持 $h$ は3または4である。
$h=3$ なら低階数分岐の証明で用いた反対点の構成を階数によらず使えることを確認する。
直線上で $w(x)=\sum_{j=1}^3 K_j(x)\otimes A_j$ と書けば、
一般の $K_j(x)$ は独立なので中央支持は一定であり、その消去点 $y$ は全 $x\in L$ と
$w(x,y)=0$ を満たす。よって $w(y)$ の中央支持は $L$ の消去二空間に含まれる。
中央1は連続純関係を与えるからちょうど2である。一般 $x$ の両端階数3と
二つの隣接六切断の直交を使うと $y$ の両端階数はどちらも3以下である。
補題\ref{lem:middle-binary-pencil}と端の下界2よりそれらは $2,3$ または $3,2$ であり、
$y$ は所望の中心である。

そこで $h=4$ とする。分岐列の各辺では中央支持の和が $\rho$ 未満である
（定理\ref{thm:low-rank-branch-exclusion}の最初の直交補の議論）。
以下 $\rho=m+5$ のときも、$L$ 上に中央支持2の中心があることを示す。
反対切断の $W_p^\perp$ はこのとき五次元であり、
直線 $x(t)$ の各中支持を、位置3の四文字を入力とする $5\times4$ の一次行列 $H(t)$ で表す。
一般階数は4である。四次小行列式の共通零因子の次数を $\delta$ とすると、
その行像の消去直線を表す原始な五つの同次式の次数は $4-\delta$。
分岐列の次の曲線 $x_2(t)$ の次数は $d_2=3-f_1$ である。
$w(-,x_2(t))$ の全ての行を $W_p$ を法として読むと、この消去直線に値を持つ
$d_2$ 次の多項式列である。全列が零なら全 $x\in L$ と全 $x_2(t)$ が隣接し、
曲線の積が辺に含まれて可動中心の排除に反する。従って原始列の次数は $d_2$ 以下で
\[
                           \delta\ge f_1+1 .
\]

他方で $H$ の階数落ちの長さを局所的に上から評価する。
先の $K(t)$ は三列の $4\times3$ 行列で、三次小行列式の局所共通零次数を $n$ とすると
全点での和は $f_1$ である。$n>0$ は順中心で、そこでの階数は2なので
局所DVRでの単因子は $1,1,t^n$（今は $t$ を一様化変数とする）。
$H(t)$ は、$A_j$ が与える定数の四次元部分空間に $K(t)\otimes 1$ を制限したものである。
実際 $A_j$ による写像は一般の中央支持4より単射で、
$W_p^\perp$ への定数包含は局所単因子を変えない。
この制限の各単因子の指数は $n$ 以下である。飽和像の元に $t^n$ を掛ければ
$K(t)\otimes1$ の全像に戻り、その逆像は一般体上で元の定数部分空間に属するので
DVR上でもその部分空間に戻る。ゆえに $\delta$ の各点での寄与は
$(4-\operatorname{rank}H(t))n$ 以下となる。
全ての順中心が中央支持3以上なら和は $f_1$ 以下となり、先の下界に反する。
従って $L$ 上に中央支持2の順中心がある。

中央支持2の中心がないと仮定すると、$m\ge3,h=4$ かつ $m+4<\rho$。
階数8なら $m=3$、階数9で $m=4$ の場合も上の $\rho=m+5$ に入り矛盾する。
残る階数9の全中心は $m=3$ であり、逆向きにも同じである。
\end{proof}
